\chapter{字符串基础与符号约定}
% \addcontentsline{toc}{chapter}{字符串基础与符号约定}

本节约定使用的基本符号和定义。除非特别说明，字符串下标从 $1$ 开始。

\section{基本符号}

\begin{center}
\renewcommand{\arraystretch}{1.25}
\begin{tabular}{p{0.24\textwidth}p{0.68\textwidth}}
$\Sigma$ & 字符集，通常假设为小写字母集合。一般来说可以是任意有限全序集。\\
$c\in \Sigma$ & 字符。\\
$s,t,S,T$ & 字符串，即有限个字符组成的序列。\\
$\epsilon$ & 空串，长度为 $0$。\\
$\Sigma^*$ & 由 $\Sigma$ 中的字符组成的所有字符串的集合。\\
$|s|$ & 字符串 $s$ 的长度。\\
$s[i]$ & $s$ 的第 $i$ 个字符。特别地，当 $i<0$ 时表示 $s[|s|+i+1]$，例如 $s[-1]=s[|s|]$。\\
$s[l,r]$、$s[l,r)$、$s(l,r)$ 等 & 子串，即区间内的 $s[i]$ 从左到右拼接形成的字符串。其中 $[l,r]$ 为闭区间，$[l,r)$ 不包含 $r$，$(l,r)$ 不包含两端。端点为负数时按上面的规则换算，例如 $s[-r,-l]=s[|s|-r+1,|s|-l+1]$。$l,r$ 中间的逗号有时会使用冒号替代。\\
$st$ 或 $s+t$ & 字符串拼接。\\
$s^R$ & 字符串 $s$ 的反串，即 $s^R[i]=s[|s|-i+1]$。\\
$s[:i]$ & 前缀 $s[1,i]$。\\
$s[i:]$ & 后缀 $s[i,|s|]$。\\
\end{tabular}
\end{center}

\section{基本概念}

\begin{center}
\renewcommand{\arraystretch}{1.25}
\begin{tabular}{p{0.24\textwidth}p{0.68\textwidth}}
前缀 & 若存在字符串 $u$ 使 $t=su$，则称 $s$ 是 $t$ 的前缀。\\
后缀 & 若存在字符串 $u$ 使 $t=us$，则称 $s$ 是 $t$ 的后缀。\\
子串 & 若存在字符串 $u,v$ 使 $t=usv$，则称 $s$ 是 $t$ 的子串。\\
子序列 & 从 $t$ 中删除若干字符后得到的字符串。子序列不要求连续。\\
$\mathrm{LCP}(s,t)$ & $s,t$ 的最长公共前缀长度。\\
$\mathrm{LCS}(s,t)$ & $s,t$ 的最长公共后缀长度。\\
字典序 & 从左到右比较第一处不同字符；若一者是另一者前缀，则较短者更小。\\
回文串 & 满足 $s=s^R$ 的字符串。\\
\end{tabular}
\end{center}


% \section*{本讲义中的核心符号}

% \begin{center}
% \renewcommand{\arraystretch}{1.25}
% \begin{tabular}{p{0.24\textwidth}p{0.68\textwidth}}
% border & $s$ 的既是前缀又是后缀的真子串。视情况也将空串和 $s$ 本身计入\\
% 周期 $p$ & 满足 $1\le p\le |s|$ 且 $s[i]=s[i+p]$ 对所有 $1\le i\le |s|-p$ 成立的整数。\\
% 整周期 $p$ & 周期 $p$ 且 $p\mid |s|$。\\
% $f(i)$ & KMP 中 $s[1,i]$ 的最长 border 长度，也常称 fail/next/前缀函数。\\
% $z(i)$ & Z 函数，表示 $\mathrm{LCP}(s[i,|s|],s)$。\\
% $h_i$ & Manacher 中以 $i$ 为中心的最长回文半径。\\
% $\delta(u,c)$ & 自动机在状态 $u$ 读入字符 $c$ 后的转移。\\
% \texttt{fail} / \texttt{link} & 失配链接、后缀链接或回文树链接；具体含义随结构而定。\\
% \end{tabular}
% \end{center}
